\section{The Glass Bead Game}

Lately, in comments, emails, and even in person, I have heard the complaint that Gornahoor is too ``esoteric'' and difficult to understand. Perhaps there is some truth in that based on some random comments I have heard that really miss the mark. Nevertheless, it is not deliberately obscure. There is a certain style and vocabulary that will weed out the uneducated, but the underlying material should be comprehensible. Moreover, I make use of some rhetorical devices which are in likelihood missed since proper rhetoric is no longer taught.

To make it clear, whatever else you may have heard. Our purpose is to awaken the memory of an intellectual conversion, i.e., to see the world as if Spirit is the fundamental reality. Although there is logic and words, discursive thought must end in that realization, as it is not an end in itself.

Thought is not to be used as a tribal marker or, as someone recently told me, it reveals to him ``who is shooting in the same direction.'' It does not matter to him that the arguments are weak, the facts are wrong, or that it may promote and accept moral depravity. This fellow claimed it is not right to critique those who are on the ``same side'', even when they do not adhere to the same standard. I couldn't convince him that we are on the side of the ``transcendent'', which is not a side.

With some young men, the discussion will always return to Nietzsche. Of course, I explained I am the ultimate Nietzschean since I accept that there are only perspectives and worldviews; hence, even those ``on his side'' have no absolute perspective. So rather than a useless, infinite (i.e., eternally recurring) battle of perspectives, the issue really devolves to this: who has the strength or the will to create a greater or more encompassing perspective?

This is the premise of \textbf{Herman Hesse}'s \emph{Glass Bead Game}. In the novel of that name, an elite engage in a project of synthesizing human knowledge. This requires knowledge of mathematics, music, history, science, philosophy and so on, along with the ability to grasp their inner inter-connectedness. So Gornahoor readers should be able to discern the influence of that Game. Along with those topics, we include the world's traditions as well as esoteric and Hermetic teachings.

In the novel, this was reserved for the few, although in our age every man deems himself capable of playing. Of course, that is why Plato made knowledge of maths a prerequisite for the study of the much more difficult fields of politics and metaphysics. Perhaps someday I will try to master LaTeX\footnote{\url{http://www.latex-project.org/}} so I can post an occasional piece on algebraic topology. Those ignorant or incapable of doing maths always know it beyond doubt. Unfortunately, that does not apply to politics or metaphysics, especially in our democratic and protestant age. Please refer to Maths and Politics\footnote{\url{https://gornahoor.net/?p=37}}.

That is why I enjoy sports, as there is little doubt in a box score. As C S Lewis once pointed out, the sports pages in the newspaper usually gets everything right. But a sports match always has a referee. Sadly, political and metaphysical discussions never have impartial referees to enforce the rules of logic, do fact checking, etc. Hence, in the glass bead game, we will have to rely on the honour system.

\paragraph{Rules of the Game}


To start the glass bead game, there must be a few preliminary rules.

The end of \textbf{egalitarianism}. Since perspectives can be judged on their coherence, consistency, rationality, etc., there can no equality of perspectives. I heard that someone recently advocated that egalitarians should be ``shamed''. As the proverb says, we should eat our own dog food. Hence, such anti-egalitarians need to be willing to have their own views critiqued and humiliated. That can only result in better, more accurate ``shooters''

The end of \textbf{parrotism} (mindless imitation): This is the mystique of quoting without understanding. Hence, we see long passages quoted without any added value or commentary. Even worse, certain terms and notions are taken out of context and used in completely inappropriate ways. The game requires that ideas need to be rephrased, combined with different but related ideas, etc. This is the only way to develop or demonstrate full understanding.

\paragraph{Justice in Thinking}


There are some rules on the positive side. The fundamental rule is justice, since every proper game requires fairness and openness. The strong man takes no pleasure in beating the weak and the inept. Justice requires the proper balance among the soul faculties of thought, will, and feeling. For the Game, the virtues that lead to justice will manifest as follows:

\textbf{Wisdom}: Following one's logic to its end. As Nietzsche wrote: ``Wer sagt A muss B sagen.'' Surprisingly few men are willing to accept all the logical consequences of their ideas, but that is a foul in the game. Of course, skillful players often let the readers come to their own conclusions. Not everything needs to be spelled out so clearly and there are often sound reasons for not doing so.

\textbf{Courage}: bucking the trend, not choosing sides. This is the opposite of my friend's claim that ideas define ``sides'' in some imaginary battle.

\textbf{Temperance}: detachment from emotionalism. This is too obvious.


\paragraph{Some topical ideas}


Here are some topic ideas that perhaps can be developed further. There is the so-called Dark Enlightenment, claimed by secularists who have come to realize a thing or two that have always been known. That is why we occasionally promote August Comte and Charles Maurras, just for them. Comte was a mathematician who tried to organize all the sciences, from mathematics and astronomy to politics and ethics. For that, he deserves an honorable mention as a player. Maurras followed his example. Although secularists and atheists (Maurras for most of his life), they did not try to overthrow the established order of things. Hence, their political programs were consistent with traditional notions of social order. It is a failing that their wisdom has been lost.

In my brief flirtation with disco many years ago, I learned what now seems to be the theme song of alt/new right\footnote{\url{https://www.youtube.com/watch?v=mGPlK2xHJyI}} ... at least you should recognize two out of three. Is there anything more than that message, after you cut through the persiflage?

There may be instant enlightenment for anyone who can answer the following poll question correctly. Or else there may be total confusion.

What is nondual heavy metal music?

\begin{itemize}


\item The song is played on just one chord.

\item The song is played on not two chords.

\item The song uses no chords at all.

\item The song uses just the white notes on the guitar.

\item All and none of the above.
\end{itemize}

Jung wrote the following on the ``shadow'', the dark, unconscious, part of the personality. We alluded to this in post called Privation\footnote{\url{https://gornahoor.net/?p=1815}} several years ago. If Jung says ``A'', I am sure some of you can then understand ``B'' for yourselves.

\begin{quotex}
I have frequently observed in the analysis of Americans that the inferior side of the personality, the ``shadow', is represented by a Negro or an Indian, whereas in the dream of a European it would be represented by a somewhat shady individual of his own kind. These representatives of the so-called `lower races' stand for the inferior personality component of the man.
\end{quotex}


The Shadow also appears among women who try to get involved with Tradition. They are attracted to the bizarre, the macabre, the weird, and seem to believe that is the meaning of ``Tradition''. Included in this, is an association with bondage or domination and that H P Lovecraft is somehow a ``traditionalist''. I suppose this should not be surprising given the immense worldwide popularity of the ``50 Shades of Gray''.

Speaking of women, I understand a new/alt right group, after its recent convention, is asking itself why there are so few women in the ``movement''. The real mystery is why any men would get involved in such a spiritually arid group.

Someone --- active in such circles --- wrote a few months ago that Gornahoor was the best Traditional blog on the Internet. Now he is less happy about it, and thinks there is little more here than a few interesting translations. Fickleness in a woman can be part of her charm, but an inconstant man is reprehensible and cannot be trusted.

Some of you may believe I just make things up. If so, consider this quote from Valentin Tomberg, which reflects ancient wisdom. When you consciously remember this about yourself, you have mastered the Game and no longer need to play.

\begin{quotex}
The individuality descends consciously and of his own free will to birth, into an environment where he is wanted and awaited ... Incarnation-birth presupposes conscious agreement between the will of the individuality above and the receiving will below. All incarnation-births are announced, i.e., preceded by knowledge of the individuality who is going to incarnate himself due either to direct intuition or to intuition revealing itself in a dream or, lastly, to revelation by means of a vision experienced by the future parents in full waking consciousness.
\end{quotex}



\flrightit{Posted on 2013-12-01 by Cologero}

\begin{center}* * *\end{center}

\begin{footnotesize}
\begin{sffamily}
\texttt{Ea Lord of the Deeps on 2013-12-02 at 06:53 said:}

Gornahoor is not `hard' or easy. It just is. I disagree of some of your opinions, of course. Not in this post tho. Its quite sharp.

I think non egalitarians need to be more aggresive in a non-emotional attached way. I rather quote the mother of the Sultan Muhammad after the fall of Granada : `Do not weep like a woman for what you could not defend like a man'.

I've found interesting the poll. I've voted `all and none of the above' in the sense of `sunyata'

\hfill

\texttt{scardanelli on 2013-12-02 at 12:13 said:}

Interesting... I had never considered that incarnation or the first birth might happen in a way analogous to the second birth. Or perhaps I already knew that and I just now remembered.

I've never understood the whole alt right black metal thing. The soul requires beauty as the lungs require air. As an antidote, listen to Ola Gjielo's Dark Night of the Soul... where beauty is there you will find truth also.

\hfill

\texttt{Avery on 2013-12-02 at 12:15 said:}

\emph{Speaking of women, I understand a new/alt right group, after its recent convention, is asking itself why there are so few women in the ``movement''. The real mystery is why any men would get involved in such a spiritually arid group.}

Having to discuss this question shows that they recognize they are falling short in some way, even if they can only diagnose the problem using the shallow language of pop feminism. I can and do eagerly share this blog with my female friends, as I might share John Michael Greer's blog. I would not do that with Mencius Moldbug. Perhaps I am reading Moldbug just to keep an eye on him.

While we're on the subject of the Dark Enlightenment (there is no need to put this name in quotes because it is extremely apt), I'd like to point out that the Japanese neo-reactionary blogosphere is much larger and more developed, but has produced no good fruits at all, and is on the contrary irresponsible and uncontrolled, encourages people to loudly proclaim their previously quiet feelings of rage and discontent, and proposes political solutions to spiritual problems. I have been Facebook friends with a Japanese net-rightist for a while now and it is clear she has contracted a spiritual sickness. The true traditionalists understood the shortcomings of fascism the first time around.

\hfill

\texttt{Jacob on 2013-12-02 at 15:05 said:}

Yea, there will always be people who can't grasp what others can. Of course (and especially for the people who god their way here), I sometimes wonder if they truly do not understand or intentionally misunderstand. Not that it matters, since the goal isn't to change everyone's mind, but to lead those who can grasp it to full Truth. Truth is never really effected by who believes it or not.

@Avery: this is slightly off topic, but is Japan just as bad as America and Europe? After reading Ash's blog with the Yukio Mushima quote I wondered if maybe Japan might have any significant opposition to modernism within it.

\hfill

\texttt{Ea Lord of the Deeps on 2013-12-02 at 16:39 said:}

To scardanelli ``alt right black metal thing'' . I think Its too look `edgy'. I only like Burzum, the rest of the `genre' is worthless noise. I like some `Blackgaze' too like Agalloch. But you're right.

\hfill

\texttt{Avery on 2013-12-02 at 17:49 said:}

@Jacob: Japanese society is liberal in superficial, legal ways, but it is resistant to the modern outlook at its core, and few people can explain why. If we really want to play this ``glass bead game'' on a worldwide scale this is something worth thinking about, although I didn't get enough of a response here and on my personal blog to develop a discussion. As I noted below, their equivalent of the Dark Enlightenment is nasty, abusive, and anti-social, so that's certainly not it.

The Japanese minds whom I consider traditional and respect, such as Mishima, \textit{Shiba Ryotaro}\footnote{\url{http://avery.morrow.name/blog/2013/06/the-beauty-of-unturned-stones/}}, \textit{Okawa Shumei}\footnote{\url{http://www.gornahoor.net/?tag=okawa-shumei}}, and \textit{Tomofusa Kure}\footnote{\url{http://avery.morrow.name/blog/2013/03/kure-tomofusa-on-mishima-yukio-and-devotion/}}, stayed in dialogue with leftists throughout their lives. This shouldn't be shocking; G.K. Chesterton considered himself a classical liberal, after all. These intellects generally preferred the language of Confucianism and bushido. Buddhism is virtually absent as a living tradition except in degenerate forms, and Taoism/Shinto, while definitely Traditional, almost plays a Dionysian role against mainstream society.

\hfill

\texttt{X on 2013-12-03 at 01:44 said:}

``Buddhism is virtually absent as a living tradition except in degenerate forms''

Avery, what did you exactly mean by this? Can you please explain it further?

\hfill

\texttt{Avery on 2013-12-03 at 03:15 said:}

This is a radical personal opinion, but I've stayed in several Japanese Buddhist temples and they seem to be traditional-style funeral parlors. The medieval forms remain, and that's good, but the disinterest of secular society in anything beside funerals means that I do not expect spiritual transformation from that quarter. The ``degenerate forms'' I referred to are the Nichiren movements, such as Soka Gakkai.

I am still investigating Shinto/Taoism, but Confucianism remains seriously overlooked as a method of understanding Tradition in the East. Apologies for derailing the thread.

\hfill

\texttt{Gina on 2013-12-03 at 16:19 said:}

This is my second week of reading this blog after I googled the logres different ways. I'm beginning to understand more about Traditionalism. The posts thus far are engaging on a semi-Sloterdjik level, which I didn't think was possible for blogs. It definitely makes for interesting inner dialogue.

\hfill

\end{sffamily}
\end{footnotesize}

